% =====================================================================
% BAGHEWALA WELL-TO-SURFACE DIGITAL TWIN — PROJECT REPORT + USER GUIDE
% SIH 2026 Problem Statement SIH26120 (Oil India Limited)
% ---------------------------------------------------------------------
% HOW TO COMPILE (single file, no .bib needed):
%   pdflatex Baghewala_DigitalTwin_ProjectReport.tex
%   pdflatex Baghewala_DigitalTwin_ProjectReport.tex   (2nd run for TOC)
% Requires a standard TeX Live / MiKTeX installation.
% Screenshots: create a "figures/" folder next to this file and drop in the
% PNGs listed in Appendix C; the placeholders will pick them up automatically.
% =====================================================================
\documentclass[12pt,a4paper]{report}

% ------------------------------- packages ----------------------------
\usepackage[utf8]{inputenc}
\usepackage[T1]{fontenc}
\usepackage{lmodern}
\usepackage[a4paper,left=2.8cm,right=2.5cm,top=2.5cm,bottom=2.5cm]{geometry}
\usepackage{xcolor}
\definecolor{accent}{HTML}{1F6F4A}
\definecolor{amber}{HTML}{B7792B}
\definecolor{codebg}{HTML}{F4F4F2}
\usepackage{amsmath,amssymb}
\usepackage{graphicx}
\usepackage{booktabs}
\usepackage{longtable}
\usepackage{enumitem}
\usepackage{caption}
\usepackage{subcaption}
\usepackage{listings}
\usepackage{tikz}
\usetikzlibrary{arrows.meta,positioning,shapes.geometries}
\usepackage{fancyhdr}
\usepackage[hidelinks]{hyperref}
\hypersetup{pdftitle={Baghewala Well-to-Surface Digital Twin — Project Report},
  pdfauthor={Project Team (SIH26120)}}

\lstdefinestyle{ts}{
  backgroundcolor=\color{codebg},
  basicstyle=\ttfamily\small,
  keywordstyle=\color{accent}\bfseries,
  commentstyle=\color{gray},
  stringstyle=\color{amber},
  numbers=left, numberstyle=\tiny\color{gray}, numbersep=8pt,
  breaklines=true, breakatwhitespace=true,
  frame=single, rulecolor=\color{gray!40},
  showstringspaces=false, tabsize=2
}
\lstset{style=ts}

\pagestyle{fancy}
\fancyhf{}
\fancyhead[L]{\small\textit{Baghewala Digital Twin (SIH26120)}}
\fancyhead[R]{\small\thepage}
\fancyfoot[C]{\small\textit{Project Report + Explanatory Guide}}
\renewcommand{\headrulewidth}{0.6pt}

\setlist[itemize]{leftmargin=*,itemsep=3pt}
\setlist[enumerate]{leftmargin=*,itemsep=3pt}

% ============================== document =============================
\begin{document}

% ------------------------------ title page ---------------------------
\begin{titlepage}
\centering
{\large\bfseries SMART INDIA HACKATHON 2026\par}
\vspace{4pt}
{\normalsize Problem Statement SIH26120 \textbar{} Oil India Limited\par}
\vspace{14pt}
\rule{\textwidth}{1.2pt}\par\vspace{10pt}
{\Huge\bfseries Digital Twin for Well-to-Surface\\[4pt]
Optimization of Cyclic Steam Stimulation\\[4pt]
(CSS) and Sucker Rod Pump (SRP)\\[4pt]
Operations for Heavy Oil Wells\\[4pt]
of Baghewala Field\par}
\vspace{10pt}
\rule{\textwidth}{1.2pt}\par
\vspace{14pt}
{\Large Project Report \textbar{} Explanatory Guide \textbar{} User Manual\par}
\vspace{22pt}
\begin{tabular}{rl}
\textit{Field:} & Baghewala PML, Bikaner--Nagaur Basin, Rajasthan \\
\textit{Reservoir:} & Jodhpur Sandstone (\textasciitilde1,110\,m, 14--19$^\circ$API heavy oil) \\
\textit{Stack:} & React 18 + TypeScript + Vite + Leaflet/OSM + Recharts \\
\textit{Deployment:} & Render Web Service (Node) serving the production bundle \\
\end{tabular}
\vspace{24pt}

\begin{center}
\begin{tabular}{p{5.5cm}p{7.5cm}}
\toprule
\textbf{Role} & \textbf{Name} \\
\midrule
Team Leader & \rule{6.5cm}{0.4pt} \\
Member 2 & \rule{6.5cm}{0.4pt} \\
Member 3 & \rule{6.5cm}{0.4pt} \\
Member 4 & \rule{6.5cm}{0.4pt} \\
Member 5 & \rule{6.5cm}{0.4pt} \\
Member 6 & \rule{6.5cm}{0.4pt} \\
Faculty Mentor & \rule{6.5cm}{0.4pt} \\
\bottomrule
\end{tabular}
\end{center}
\vspace{18pt}
{\large Institution: \rule{8cm}{0.4pt}\par}
\vspace{6pt}
{\large September 2026\par}
\end{titlepage}

% ---------------------------- certificate ----------------------------
\chapter*{Certificate}
\addcontentsline{toc}{chapter}{Certificate}
This is to certify that the project titled \textbf{``Digital Twin for
Well-to-Surface Optimization of Cyclic Steam Stimulation (CSS) and Sucker Rod
Pump (SRP) Operations for Heavy Oil Wells of Baghewala Field''} is a bonafide
record of work carried out by the team listed on the title page, in partial
fulfilment of the requirements of Smart India Hackathon 2026, Problem Statement
SIH26120 floated by Oil India Limited. The work is original, and all sources of
field data and engineering correlations used herein have been duly cited.

\vspace{30pt}
\begin{tabular}{p{7cm}p{7cm}}
\rule{5.5cm}{0.4pt} & \rule{5.5cm}{0.4pt} \\
Faculty Mentor (Signature) & Head of Department (Signature) \\
 & \\
\rule{5.5cm}{0.4pt} & \rule{5.5cm}{0.4pt} \\
External Examiner & Date / Seal \\
\end{tabular}

% ---------------------------- declaration ----------------------------
\chapter*{Declaration}
\addcontentsline{toc}{chapter}{Declaration}
We declare that (i) this report is our own work except where acknowledged;
(ii) field figures (depths, API gravities, viscosities, well counts, production
records) are taken from the published sources cited in Chapter~3 and reproduced
only as calibration anchors; (iii) every computed quantity in the software
(rate, SOR, temperature, viscosity, fillage, load, margin) is \emph{solved} from
the coupled engineering chain of Chapter~4 --- no production number is invented
by a random generator; and (iv) the well layout is a representative 12-pad grid
inside the real Baghewala PML and is explicitly labelled as such wherever shown.

\vspace{24pt}
\begin{tabular}{ll}
Date: \rule{4cm}{0.4pt} & Place: \rule{5cm}{0.4pt} \\
\end{tabular}
\vspace{12pt}

\noindent Signatures of team members:\quad
\rule{3cm}{0.4pt}\quad \rule{3cm}{0.4pt}\quad \rule{3cm}{0.4pt}\quad
\rule{3cm}{0.4pt}\quad \rule{3cm}{0.4pt}\quad \rule{3cm}{0.4pt}

% --------------------------- acknowledgement -------------------------
\chapter*{Acknowledgement}
\addcontentsline{toc}{chapter}{Acknowledgement}
We thank Oil India Limited for floating a problem statement grounded in a real
operating field; the authors of the SPE, AIME and API publications whose
correlations power our model; and the OpenStreetMap contributors whose open map
data lets our field view show real geography instead of rendered backdrops.

% ------------------------------ abstract ----------------------------
\chapter*{Abstract}
\addcontentsline{toc}{chapter}{Abstract}
Baghewala field (Rajasthan) produces 14--19$^\circ$API heavy oil
(8\,000--15\,000\,cP at 50$^\circ$C) from the Jodhpur Sandstone at
\textasciitilde1,110\,m, virgin temperature 46--48$^\circ$C and low reservoir
pressure. Recovery needs Cyclic Steam Stimulation (CSS) plus sucker-rod pumping
(SRP), but the two are today optimised separately: steam slugs from historical
practice, pump speed adjusted reactively. As the rock cools, viscosity climbs,
pump fillage collapses, rods float and fail, and steam-oil ratio (SOR) rises.

We built an AI-enabled, well-to-surface digital twin as a web application that
\emph{couples} both halves in one deterministic chain:
\begin{center}
CSS slug $\rightarrow$ Marx--Langenheim heated radius $\rightarrow$ soak
retention $\rightarrow$ $\mu(T)$ $\rightarrow$ Vogel IPR inflow
$\rightarrow$ API RP\,11L displacement $\rightarrow$ fillage
$\rightarrow$ rod loads $\rightarrow$ Goodman fatigue + rod-float margin
$\rightarrow$ SOR / INR operating margin.
\end{center}
A grid optimiser over steam $\times$ soak $\times$ SPM ranks every operating
node on rupees per day (oil revenue minus steam, power and a risk penalty) and
presents the Pareto knee as the recommended window. Twelve representative pads
with true WGS84 coordinates sit on live OpenStreetMap tiles; per-well rates
(15--45\,BOPD) scale to the published 1,202\,BOPD field record (Apr-2026).
A 15-stop in-app guided tour narrates the system for evaluators, and an
Engineering Basis page traces every number to its public source or equation.

\textbf{Keywords:} digital twin, CSS, sucker-rod pump, heavy oil, Baghewala,
Vogel IPR, API RP\,11L, Marx--Langenheim, SOR, OpenStreetMap.

\tableofcontents
\listoffigures
\listoftables

% =====================================================================
\chapter{Introduction}
% =====================================================================
\section{India's heavy-oil challenge}
Conventional light-oil reservoirs dominate Indian production, but heavy-oil
accumulations such as Baghewala represent a large, hard-to-produce resource.
Heavy oil is defined here pragmatically: dead-oil viscosity in the thousands of
centipoise at surface conditions, API gravity below about 20$^\circ$, and
negligible natural flow energy. Primary recovery is poor; sustained production
demands \emph{thermal} enhanced oil recovery together with \emph{artificial
lift} --- and, crucially, the two must be managed as one system.

\section{The Baghewala field}
Baghewala lies in the Bikaner--Nagaur sub-basin of the Rajasthan Basin
(Baghewala Petroleum Mining Lease, 200.26\,km$^2$). Oil India Limited's
Baghewala-1 (1991) discovered heavy oil in the Neoproterozoic--Cambrian Jodhpur
Sandstone at 1,104--1,117\,m (average pay depth \textasciitilde1,150\,m).
Key published facts used throughout this project:
\begin{itemize}
  \item Fluid: 17--19$^\circ$API per SIH26120; laboratory assays 14--17$^\circ$API
    (SPE APOG 2023, Paper 535203); viscosity 10\,000--13\,000\,cP at 50$^\circ$C
    (OIL) / 8\,000--15\,000\,cP (SPE); sulphur-rich ($>$1\,\%S), high asphaltene.
  \item Reservoir: virgin temperature 46--48$^\circ$C, low pressure, poor primary
    mobility; net pay \textasciitilde14\,m, porosity \textasciitilde19\,\%,
    permeability 300--1\,500\,mD.
  \item Development: first CSS pilot 2006, commercial production since 2017;
    35 wells drilled / 23 producing (base case) growing to 52 wells / 33
    producing with a record 1,202\,BOPD in April 2026; CSS active in 19 wells
    (+72\,\% year-on-year); thermal completions (thermal wellhead +
    vacuum-insulated tubing, VIT); conventional + hydraulic SRP units.
\end{itemize}

\section{Cyclic Steam Stimulation in brief}
A CSS cycle has three phases: (1) \textbf{injection} --- saturated steam
(typically 24--40\,bar, \textasciitilde280--340$^\circ$C) is pumped for days to
weeks (slug 300--550\,m$^3$ cold-water equivalent, CWE); (2) \textbf{soak} ---
the well is shut in (48--120\,hr) so heat conducts outward; (3)
\textbf{production} --- viscosity near the wellbore has collapsed by two orders
of magnitude and oil flows to the pump. The economic yardstick is the
steam--oil ratio, SOR $=$ (CWE steam in) / (oil out); literature band 2.5--5,
with small-slug designs targeting below 3.

\section{Sucker-rod pumping in brief}
A surface beam unit reciprocates a rod string driving a downhole plunger pump
(here 1.75-in.\ plunger at 1,075\,m, 7/8-in.\ API Grade-D rods, 2$\tfrac78$-in.\
tubing, 22\,kW motors via VFD). Operating point: strokes per minute
(3.0--6.5\,SPM), stroke length (64--86\,in), VFD frequency (35--55\,Hz). Pump
displacement follows API RP\,11L:
\begin{equation}
PD = 0.1166\, D^{2}\, S_{p}\, N \quad \text{[BPD; $D,S_{p}$ in inches, $N$ in SPM]}.
\label{eq:pd}
\end{equation}
In heavy oil the rods can \emph{float}: on the downstroke, viscous drag plus
fluid load can exceed the buoyed rod weight, so rods fall slowly and then slam
(impact loading), unsettling pumps and parting rods.

\section{Why separate optimisation fails}
The coupling chain (Figure~\ref{fig:chain}) shows the trap: raising SPM without
re-heating the near-wellbore merely churns the same viscous fluid harder ---
fillage falls, rod-float margin collapses, Goodman fatigue ratio climbs, and SOR
barely moves. Thermal and lift decisions share one state vector (pump fillage);
optimising them apart guarantees steam waste on one side and iron breakage on
the other.

\section{The digital-twin concept}
A digital twin is a live, quantitative mirror of a physical asset: field state
in, prediction and optimisation out, with every recommendation traceable to
physics or data. Our twin mirrors each well from reservoir to sales point
(``well-to-surface''), recomputing the full chain whenever any control moves.

\section{Organisation of this report}
Chapter~2 states the problem and objectives; Chapter~3 documents field data and
calibration; Chapter~4 derives the engineering models; Chapter~5 presents
system architecture; Chapter~6 the implementation; Chapter~7 is the explanatory
user guide; Chapter~8 deployment; Chapter~9 results and validation;
Chapter~10 limitations and future work; followed by references and appendices
(judge Q\&A, glossary, screenshot guide).

% =====================================================================
\chapter{Problem Statement and Objectives}
% =====================================================================
\section{SIH26120 as floated}
SIH 2026 Problem Statement SIH26120 (Oil India Limited), titled as on our
cover, observes that Baghewala's CSS cycle design and SRP operation are
optimised separately from historical experience; that post-steam cooling raises
viscosity, cutting pump efficiency and causing rod floating, impact loading,
pump unsetting and rod failures; and it asks for an AI-enabled well-to-surface
digital twin giving real-time monitoring, prediction and optimisation.

\section{Operating challenges addressed}
\begin{enumerate}
  \item CSS parameters (steam volume, injection pressure, soak time,
    production cut-off) set from historical practice, not optimisation.
  \item SRP parameters (stroke length, SPM, VFD) adjusted manually and
    reactively.
  \item Heavy crude causes rod floating, impact loading, pump unsetting, rod
    failures and heavy maintenance.
  \item Reservoir, wellbore and lift performance are never optimised together.
  \item No predictive analytics: high SOR, high energy per barrel, low recovery.
\end{enumerate}

\section{Objectives}
\begin{enumerate}
  \item Optimise CSS cycle parameters per well and per cycle.
  \item Predict heating, cooling and production trajectories.
  \item Continuously recommend SRP operating points (SPM, stroke, VFD) from
    live well state.
  \item Detect rod floating and minimise impact loading with force-based
    criteria.
  \item Improve pump efficiency and equipment reliability.
  \item Minimise steam and energy per barrel (operating margin in INR/day).
\end{enumerate}

\section{Expected benefits}
Higher recovery; lower SOR; lower energy per barrel; fewer rod failures and
unset pumps; longer equipment life; and, above all, predictive, data-driven
decisions an operations engineer can audit.

\section{Scope and assumptions}
The twin models 12 representative pads inside the real PML (the 52-well pattern
compressed for demonstration; totals scale linearly). It is calibrated to
\emph{published} data, not live SCADA --- the Sensor Integrity page
demonstrates the drift detect $\rightarrow$ exclude $\rightarrow$ recalibrate
architecture a live feed would use. Economic constants: oil USD\,78/bbl at
INR\,83.5, steam INR\,1\,450/m$^3$ CWE, power INR\,8.5/kWh.

% =====================================================================
\chapter{Field Data and Calibration}
% =====================================================================
\section{Published sources}
Every field number in the twin is traceable to Table~\ref{tab:sources}. Where
sources give ranges, the code stores the range \emph{and} the single
calibration value (see \texttt{src/engineering/fieldConstants.ts}), so any
evaluator can audit each assumption.

\begin{longtable}{p{3.2cm}p{6.2cm}p{5.2cm}}
\caption{Field-data provenance.\label{tab:sources}} \\ \toprule
\textbf{Claim used} & \textbf{Public source} & \textbf{Use in twin} \\ \midrule
\endfirsthead
\toprule \textbf{Claim used} & \textbf{Public source} & \textbf{Use in twin} \\ \midrule
\endhead
PML 200--200.26\,km$^2$; 35/23 then 52--56/33--34 wells & Oil India Ltd,
Rajasthan Fields page & Block size, well counts, status \\[2pt]
Pay \textasciitilde1,150\,m; 10--13k\,cP @50$^\circ$C; VIT + SRP; $>$600\,BOPD &
OIL Rajasthan Fields & Depth, $\mu$ anchor, completion, base rate \\[2pt]
API 14--17; 8--15k\,cP; pay \textasciitilde1,100\,m &
Singh et al., SPE APOG 2023 (535203) & PVT band, depth cross-check \\[2pt]
Bilara + Jodhpur heavy oil; Bikaner--Nagaur &
DGH National Data Repository & Stratigraphy, basin setting \\[2pt]
Baghewala-1 (1991), 17.6$^\circ$API @1,104--1,117\,m &
Bhat et al.\ via Springer J.\ Petrol.\ Expl.\ Prod.\ Technol.\ 2021 &
Discovery interval, perfs \\[2pt]
1,202\,BOPD record; CSS in 19 wells (+72\,\%) &
Times of India / News18 / ET Energy, Apr-2026 & Field-total calibration \\[2pt]
17--19$^\circ$API; $T_{res}$ 46--48$^\circ$C; scope &
SIH26120 problem statement & Framing, virgin $T$, scope \\
\bottomrule
\end{longtable}

\section{Model constants}
Table~\ref{tab:constants} summarises the working constants. Per-well API spans
14--19 (the union of SIH and SPE figures); pressure spans 70--95\,bar;
permeability spans 300--1\,500\,mD.

\begin{table}[h]
\centering
\caption{Working constants (representative calibration values).\label{tab:constants}}
\begin{tabular}{lll}
\toprule
\textbf{Quantity} & \textbf{Value} & \textbf{Remarks} \\ \midrule
Pay depth / perfs & 1,110\,m / 1,104--1,117\,m & Discovery interval \\
Virgin $T$ / $P$ & 47$^\circ$C / 82\,bar & Ranges 46--48$^\circ$C, 70--95\,bar \\
Net pay / $\phi$ / $k$ & 14\,m / 0.19 / 850\,mD & Fluvio-deltaic Jodhpur sand \\
$\mu_{50}$ anchor & 11,000\,cP @50$^\circ$C & Band 8,000--15,000\,cP \\
Steam slug / soak & 400\,m$^3$ CWE / 72\,hr & Ranges 300--550, 48--120 \\
Injection $P$ / quality & 24--40\,bar / 0.8 & Saturated steam \\
Pump / rods & 1.75-in.\ @1,075\,m; Grade-D 7/8-in. & VIT completion \\
SPM / stroke / VFD & 3.0--6.5 / 64--86\,in / 35--55\,Hz & 22\,kW motors \\
SOR target & $<3$ (band 2.5--5) & Small-slug design \\
\bottomrule
\end{tabular}
\end{table}

\section{Calibration checkpoints}
The model must reproduce, without fitting: (i) $\mu$ inside 8--15k\,cP at
50$^\circ$C for every API; (ii) 15--45\,BOPD per well, scaling to
\textasciitilde1,202\,BOPD over 33 producers; (iii) cycle SOR in 1--4;
(iv) map pads at true coordinates on real tiles. Chapter~9 verifies each.

% =====================================================================
\chapter{Engineering Models (Theory)}
% =====================================================================
This chapter derives, one by one, the correlations implemented in
\texttt{src/engineering/correlations.ts}. Numbered equations match code
comments so an evaluator can walk from formula to function.

\section{Dead-oil viscosity versus temperature}
Heavy-oil viscosity follows an Arrhenius-type collapse (Mehrotra--Svrcek
double-log form, used here in exponential form calibrated to OIL assays):
\begin{equation}
\mu(T)=\mu_{50}\,\exp[-b\,(T-50)], \qquad
b = 0.068 - 0.0016\,(\mathrm{API}-14)\;\;[/^\circ\mathrm{C}],
\label{eq:mu}
\end{equation}
with the 50$^\circ$C anchor mapped linearly across the assay band,
\begin{equation}
\mu_{50} = 8000 + \frac{19-\mathrm{API}}{19-14}\times7000
\quad\text{[cP; API\,19 $\rightarrow$ 8k, API\,14 $\rightarrow$ 15k]}.
\label{eq:mu50}
\end{equation}
Checks: $\mu(47^\circ\mathrm{C})\approx1.23\,\mu_{50}\approx13{,}500$\,cP;
$\mu(150^\circ\mathrm{C})\approx12$\,cP --- the two-order collapse that makes
CSS work --- and roughly $10\times$ per 35$^\circ$C, the heavy-oil rule of
thumb. The inverse $T(\mu)=50-\ln(\mu/\mu_{50})/b$ sizes re-steam triggers.

\section{Saturated-steam tables}
Injection pressure maps to steam temperature through an Antoine-form fit of the
IAPWS-IF97 saturation line (15--60\,bar, $\pm1.5^\circ$C), letting anyone
cross-check $P\rightarrow T$ without steam tables:
\begin{equation}
T_{sat}=\frac{3816.44}{11.6703-\ln p}-(-46.13)-273.15,
\qquad p\ \mathrm{in\ kPa},
\label{eq:tsat}
\end{equation}
e.g.\ \textasciitilde224$^\circ$C at 25\,bar. Dry-steam enthalpy is linearised
as $h_s \approx 2796-0.62\,(P-30)$\,kJ/kg over 20--40\,bar; feedwater enthalpy
$h_w=4.19\,T_r$.

\section{Marx--Langenheim heated zone}
For a radial steam soak (Marx \& Langenheim, Trans.\ AIME 216, 1959), injected
heat $Q_{inj}=V_s\rho\,(h_s-h_w)$ splits into rock heating and overburden loss:
\begin{equation}
Q_{inj}=\frac{M\,h\,A_r\,(T_s-T_r)}{G(t_D)}+Q_{loss},\qquad
t_D=\frac{4\,k_{ob}\,t}{M\,h^{2}},
\label{eq:ml}
\end{equation}
with $M\approx2.35$\,MJ/m$^3\!\cdot\!$K (Jodhpur sand), $h=14$\,m,
$k_{ob}\approx2.0$\,W/m$\cdot$K, $\alpha\approx0.8\times10^{-6}$\,m$^2$/s, and
$G(t_D)$ the Marx--Langenheim conduction function (Pad\'e-style approximation
in code). Heated radius $r_h=\sqrt{A_r/\pi}$ (clamped to the physical
700--26\,000\,m$^2$ area window) and heat-loss fraction 0.15--0.55 are reported
per slug --- the CSS page's ``heat audit''.

\section{Soak retention (Boberg--Lantz character)}
After injection stops, conduction spreads heat while the near-wellbore cools
(Boberg \& Lantz, JPT 1966, SPE\,1578-PA). Soak efficiency --- the fraction of
useful heat retained in the drainage radius --- is modelled as
\begin{equation}
\eta_{soak}\approx0.97-\mathrm{loss}(t_{soak},r_h),\qquad
T_{hz}=T_r+(T_s-T_r)\times0.52\,\eta_{soak}\times\mathrm{intensity},
\label{eq:soak}
\end{equation}
where intensity scales with slug size relative to 420\,m$^3$. Longer soak
raises retention with \emph{zero} extra steam: the cheapest barrel in the
system, and the reason the optimiser often prefers $+12$\,hr of soak over
$+25$\,m$^3$ of steam.

\section{Vogel inflow with heated mobility}
Single-phase radial productivity scaled by heated mobility $k/\mu(T)$:
\begin{equation}
J=\frac{0.00708\,k\,h\,F_p}{\mu\,B\,\ln(r_e/r_w)}
\quad\text{[bbl/d/psi]},
\label{eq:j}
\end{equation}
with $F_p=0.42$ (relative permeability to heavy oil $\times$ skin / partial
penetration), history-matched so Vogel reproduces 15--45\,BOPD/well.
Below $0.6\,p_r$ the Vogel (1968) correction applies:
\begin{equation}
q=q_b+(q_{max}-q_b)\Bigl[1-0.2\frac{p_{wf}}{p_r}
-0.8\Bigl(\frac{p_{wf}}{p_r}\Bigr)^{2}\Bigr],
\label{eq:vogel}
\end{equation}
else $q=J(p_r-p_{wf})$. Flowing pressure follows pump submergence via
$p_{wf}=p_r(1-d)$, $d=0.42+(100-F\%)/220$ clamped to $[0.3,0.72]$.

\section{API RP\,11L rod-pump mechanics}
Displacement is Eq.~\ref{eq:pd}. Peak/minimum polished-rod loads combine buoyed
rod weight, fluid load $F_o=0.433\,\gamma D^{2}H\times0.785$, Tak\'acs-type
viscous drag $F_{drag}\propto\mu^{0.35}S_pN$, and a speed-squared dynamic term.
Polished-rod power is computed from first principles --- hydraulic lift
$\rho gQH$ at 55\,\% efficiency plus rod-friction work plus 0.5\,kW drive
losses (a 100\,BPD lift from 1,075\,m is \textasciitilde6\,kW shaft: checkable
by hand).

\section{The nodal coupling: fillage}
The twin's single coupling variable is pump fillage,
\begin{equation}
f=\frac{q_{liq}}{PD},\qquad
q_{liq}=\frac{q_{oil}}{1-WC}\;\; \mathrm{clamped\ to\ }[0.18,0.98],
\label{eq:fillage}
\end{equation}
solved by 3-pass fixed-point iteration with the drawdown relation
(fillage $\rightarrow$ $p_{wf}$ $\rightarrow$ inflow $\rightarrow$ fillage).
Lifted oil $=PD\times\eta_{vol}\times(1-0.12\,WC)$, capped at 60\,BOPD.

\section{Goodman fatigue check}
With $S_e=35$\,ksi, $S_y=115$\,ksi for Grade-D steel and a 2.2 coupling
stress-concentration factor (API RP\,11BR --- pins, not bodies, fail):
\begin{equation}
\frac{\sigma_a}{S_e}+\frac{\sigma_m}{S_y}\;\le\;0.80,
\label{eq:goodman}
\end{equation}
where $\sigma_a$ carries the SCF. Operating bands: $<0.55$ normal,
$0.55$--$0.80$ watch, $\ge0.80$ endurance exceedance.

\section{Rod-float criterion}
Rods float when net downstroke force vanishes: buoyed rod weight versus fluid
plus viscous resistance. Margin in lbf; healthy $>900$, floating $<250$.
Risk score $=100\,[1-\mathrm{margin}/4500]$ clamped to 2--98. Impact-loading
risk blends SPM overspeed, Goodman exceedance and $(1-f)$ starvation.

\section{Dynamometer cards}
The 97-point surface card synthesises Everitt--Jennings (SPE\,18189)
character: fluid-pound shoulder position encodes fillage ($<$\,\textasciitilde92\,\%
shows a shoulder), viscous hysteresis fattens the loop with $\log\mu$, gas
rounds the compression corner, travelling-valve leak tapers the downstroke.
The pump card follows by static load transfer (surface minus buoyed rod
weight, $\times0.92$).

\section{Decline, cooling and economics}
Forecasts use Arps exponential decline $q(t)=q_0e^{-Dt}$ with Ramey-type
cooling $T(t)=T_v+(T_0-T_v)e^{-t/\tau}$ ($\tau\approx55$\,d for VIT at
1,100\,m) and P10--P90 bands of $\pm(5\%+0.5\%/\mathrm{day})$.
Cycle SOR $=$ CWE steam / oil in the 60-day flush window. The optimiser ranks
grid nodes on
\begin{equation}
\text{margin}=q\,p_{oil}-\frac{V_s c_s}{120}-P\,24\,c_e
-(R_{float}+R_{impact})\times55\quad\text{[INR/day]},
\label{eq:margin}
\end{equation}
with steam amortised per producing day so slugs and rates compare fairly.

% =====================================================================
\chapter{System Architecture and Design}
% =====================================================================
\section{Coupled data flow}
\begin{figure}[h]
\centering
\begin{tikzpicture}[node distance=1.15cm,
  box/.style={rectangle,draw=accent!80!black,thick,rounded corners=3pt,
    minimum width=2.6cm,minimum height=0.85cm,align=center,font=\small},
  arr/.style={-{Stealth[length=2.5mm]},thick,accent!80!black}]
  \node[box] (css) {CSS slug\\$V_s,P_{inj},t_{soak}$};
  \node[box,right=of css] (ml) {Marx--Langenheim\\$r_h,Q,GJ$};
  \node[box,right=of ml] (mu) {$\mu(T)$\\Arrhenius};
  \node[box,right=of mu] (ipr) {Vogel IPR\\$q_{oil}$};
  \node[box,below=of ipr] (fill) {Fillage\\$f=q_{liq}/PD$};
  \node[box,left=of fill] (srp) {API RP\,11L\\$PD,PPRL$};
  \node[box,left=of srp] (mech) {Goodman\\+ float margin};
  \node[box,left=of mech] (econ) {SOR / INR\\margin};
  \draw[arr] (css)--(ml); \draw[arr] (ml)--(mu); \draw[arr] (mu)--(ipr);
  \draw[arr] (ipr)--(fill);
  \node[box,above=of srp] (sp) {SRP point\\$SPM,S_p,VFD$};
  \draw[arr] (sp)--(srp); \draw[arr] (srp)--(fill);
  \draw[arr] (fill)--(mech); \draw[arr] (mech)--(econ);
\end{tikzpicture}
\caption{One-directional coupled solve: no circularity; SRP and CSS meet only
at fillage. Implemented by \texttt{solveCoupledChain()}.\label{fig:chain}}
\end{figure}

\section{Module map}
\begin{table}[h]
\centering
\caption{Repository modules (matching \texttt{src/}).\label{tab:modules}}
\begin{tabular}{lp{9.5cm}}
\toprule
\textbf{Module} & \textbf{Responsibility} \\ \midrule
\texttt{engineering/} & Sourced constants + correlations (audit trail) \\
\texttt{simulation/} & Coupled solver + INR grid optimiser \\
\texttt{data/} & 12-pad WGS84 inventory, forecast/alert/sensor providers \\
\texttt{components/} & Shared UI + live OSM \texttt{FieldMap} + tour overlay \\
\texttt{pages/} & 11 routed views (Table~\ref{tab:pages}) \\
\texttt{tour/} & 15-step judge walkthrough engine \\
\texttt{types/, utils/} & Domain types, deterministic RNG helpers \\
\texttt{styles/} & Operations-console dark theme \\
\bottomrule
\end{tabular}
\end{table}

\section{Technology stack}
\begin{table}[h]
\centering
\caption{Stack and why each piece was chosen.\label{tab:stack}}
\begin{tabular}{llp{7cm}}
\toprule
\textbf{Layer} & \textbf{Choice} & \textbf{Rationale} \\ \midrule
UI & React 18 + TypeScript & Typed domain model; component reuse \\
Build & Vite 5 (\texttt{tsc -b \&\& vite build}) & Fast, static-bundle output \\
Map & Leaflet 1.9 + react-leaflet 4.2, OSM tiles & Open-source true geography (ODbL) \\
Charts & Recharts 2.12 & IPR, sweeps, cards, Pareto, forecasts \\
Routing & react-router-dom 6 & Deep-linkable wells (\texttt{/well/BGW-07/css}) \\
Hosting & Render Web Service (Node) & Serves \texttt{dist/} via \texttt{vite preview} \\
\bottomrule
\end{tabular}
\end{table}

\section{Coupling algorithm}
\texttt{solveCoupledChain(plan, point, $T_r,p_r$, API, WC, $k$)}:
(1) Marx--Langenheim $\rightarrow$ $r_h$, GJ, loss; (2) soak efficiency
$\rightarrow$ heated-zone $T$; (3) $\mu(T)$; (4) 3-pass nodal loop
(fillage $\rightarrow$ drawdown $\rightarrow$ Vogel inflow $\rightarrow$
RP\,11L displacement $\rightarrow$ fillage); (5) lifted rate, PRHP, Goodman,
float margin/risks; (6) SOR, cycle oil, steam MWh. Legacy wrappers
(\texttt{predictCssPlan}, \texttt{predictSrpPoint}) keep page code stable.

% =====================================================================
\chapter{Implementation}
% =====================================================================
\section{Engineering layer (the audit trail)}
\texttt{fieldConstants.ts} stores FIELD / RESERVOIR / FLUID / CSS / SRP /
PRODUCTION\_BENCHMARKS plus IAPWS fits. \texttt{correlations.ts} implements
Eqs.~\ref{eq:mu}--\ref{eq:margin} with reference tags
\texttt{[ML-1959]}, \texttt{[BL-1966]}, \texttt{[Vogel-68]},
\texttt{[API-11L]}, \texttt{[API-11E]}, \texttt{[EJ-1992]},
\texttt{[Mehrotra]}, \texttt{[Arps-45]}. Example:
\begin{lstlisting}[language=C]
// Per-well anchor: API 19 -> 8000 cP, API 14 -> 15000 cP
export function mu50ForApi(apiGravity: number): number {
  const [lo, hi] = FLUID.mu50RangeCp;
  const t = clamp((19 - apiGravity) / (19 - 14), 0, 1);
  return lo + t * (hi - lo);
}
export function viscosityCpAtTemp(tempC, apiGravity): number {
  const mu50 = mu50ForApi(apiGravity);
  const b = 0.068 - 0.0016 * (apiGravity - 14);
  return clamp(mu50 * Math.exp(-b * (tempC - 50)), 8, 22000);
}
\end{lstlisting}

\section{Simulation layer}
\texttt{physics.ts} exposes the coupled solver, the CSS/SRP predictors, and
\texttt{optimizeWell()}, which enumerates steam $\times$ soak $\times$ SPM
nodes and returns the best-margin node plus the ranked frontier for the Pareto
chart and top-8 table.

\section{Data providers}
\texttt{wellProvider.ts} builds the 12 deterministic pads (centre
27.58$^\circ$N, 72.82$^\circ$E; \textasciitilde5$\times$3\,km surveyed-style
grid; BGW-07: cycle 5, API 15.1, 680\,mD) by running each pad's slug and pump
point through the solver --- so inventory numbers are \emph{outputs}, not
inputs. Forecast, alert (thresholds: float 900\,lbf, Goodman 0.80, fillage
60\,\%, SOR 4.0, water cut 65\,\%) and sensor (6 channels $\times$ 12 pads;
scripted BGW-07 thermocouple drift $+8.4^\circ$C, residual gate $\pm1.5^\circ$C)
providers follow the same pattern.

\section{Views}
\begin{table}[h]
\centering
\caption{Routed pages and their purpose.\label{tab:pages}}
\begin{tabular}{lp{9.8cm}}
\toprule
\textbf{Route} & \textbf{Purpose} \\ \midrule
\texttt{/} Mission Brief & Hook, coupling chain, calibration badges, demo paths \\
\texttt{/field} & Live OSM map, KPIs scaled to 1,202\,BOPD, 14-day trend, pad cards \\
\texttt{/well/:id} & Schematic, state vector, IPR + $\mu(T)$ charts, heat audit \\
\texttt{/well/:id/css} & Slug sliders, heat audit, steam/soak sweep charts \\
\texttt{/well/:id/srp} & SPM/stroke/VFD sliders, surface + pump cards, SPM cliff sweep \\
\texttt{/well/:id/optimizer} & Recommended window, Pareto frontier, top-8, ``why'' panel \\
\texttt{/well/:id/forecast} & 7/14/30-day P10--P90 rate, BHT/$\mu$, fillage/water-cut \\
\texttt{/well/:id/scenarios} & Baseline-vs-scenario lab with live margin $\Delta$ \\
\texttt{/risk} & Threshold table + ranked mechanism/action alert cards \\
\texttt{/sensors} & SCADA channel map, drift detect $\rightarrow$ exclude $\rightarrow$ recalibrate \\
\texttt{/references} & Field + model provenance tables (the audit page) \\
\bottomrule
\end{tabular}
\end{table}

\section{Guided-tour subsystem}
\texttt{tour/steps.ts} (15 stops: route + \texttt{data-tour} target + narration
+ optional hint), \texttt{tour/TourContext.tsx} (active step, autoplay
\textasciitilde14\,s, arrow-key/Esc handling, auto-navigation), and
\texttt{components/TourOverlay.tsx} (pointer-transparent dim layer, amber
spotlight ring measured from the target's bounding box, narration card with
progress dots). Entry points: sidebar, topbar and hero ``Guided Tour'' buttons.

% =====================================================================
\chapter{Explanatory User Guide}
% =====================================================================
\section{Setup}
Prerequisites: Node.js 18+ and npm. Internet access is required only for OSM
map tiles; all computation is local.
\begin{lstlisting}
npm install
npm run dev        # open the printed URL, typically http://localhost:5173
npm run build      # type-checks (tsc -b) and emits dist/
npm run preview    # serves dist/ locally
npm start          # production serve: vite preview --host 0.0.0.0 --port ${PORT:-4173}
\end{lstlisting}

\section{Page-by-page walkthrough}
\paragraph{Mission Brief (/).} Read the hook, then the coupling chain, then the
calibration badges. Press \textbf{Guided Tour} here for the narrated version.

\paragraph{Field Overview (/field).} The map proves geography: zoom to
Baghewala and compare with any atlas. KPIs show the 12-pad total beside the
$\times33$-well scaling to the 1,202\,BOPD record. Click BGW-07's marker
$\rightarrow$ popup $\rightarrow$ twin.

\paragraph{Well Twin (/well/BGW-07).} The schematic anchors depths; the IPR
chart shows heated deliverability; the log-scale $\mu(T)$ chart shows how heat
buys it. The state vector lists every live variable with units and tones.

\paragraph{CSS Optimization.} Drag steam 410 $\rightarrow$ 500\,m$^3$: radius
grows $\propto\sqrt{Q}$, SOR $\propto Q$. Reset, then add $+12$\,hr soak:
retention rises for free. The steam/soak sweeps make diminishing returns
visible; the heat audit accounts every GJ.

\paragraph{SRP Optimization.} Drag SPM to 6.6: the card shoulder slides as
fillage collapses; the sweep shows production plateauing while Goodman climbs.
Return to 5.4. Read Goodman, float margin (lbf) and drag (lbf) as a force
balance.

\paragraph{Surface Optimizer.} Hover frontier bubbles (margin in INR/day per
bubble); read the top-8 table; verify each ``why'' bullet against its numbers.
Try arguing: change one control in What-If Lab and watch margin $\Delta$.

\paragraph{Forecast.} Switch 7/14/30-day horizons; note the widening band.
Find the BHT $\approx$ 60$^\circ$C crossing --- that is the re-steam trigger,
visible weeks early.

\paragraph{What-If Lab.} The header margin $\Delta$ settles every debate:
green earns, red loses, and the baseline/scenario panels show the physical
reason (radius, viscosity, fillage, Goodman).

\paragraph{Risk \& Alerts.} Read thresholds first, then cards: mechanism in
units, action in operations language. Disagree with a threshold, never guess
one.

\paragraph{Sensor Integrity.} Note the $+8.4^\circ$C residual against the
$\pm1.5^\circ$C gate; press recalibrate and watch confidence restore. While
flagged, recommendations use model BHT plus five corroborating channels.

\paragraph{Engineering Basis.} Trace any number here: field claims to public
sources, equations to references. End every demo on this page.

\section{Guided tour (evaluator mode)}
Fifteen stops mirror Section~7.2 with on-screen narration, hints (``try: push
SPM to 6.6''), autoplay, dots for random access, and Esc to exit. Use it for
live judging; use \texttt{VIDEO\_SCRIPT.md} for the recorded video.

\section{Demo scripts}
\textbf{3-minute judge path:} Mission Brief $\rightarrow$ Field Overview
$\rightarrow$ BGW-07 $\rightarrow$ CSS slider move $\rightarrow$ SRP slider
move $\rightarrow$ Optimizer knee $\rightarrow$ Risk threshold
$\rightarrow$ Sensor recalibrate $\rightarrow$ Engineering Basis.
\textbf{8--9 minute video:} follow \texttt{VIDEO\_SCRIPT.md} scenes 1--9
verbatim (hook, problem, map, twin, CSS, SRP, optimiser, risk+sensors, audit
close), \textasciitilde1,230 words at \textasciitilde140\,wpm.

\section{Troubleshooting}
Map grey/empty: no internet (tiles are live). Port busy: change \texttt{--port}.
Charts blank after code edits: rerun build (type errors fail \texttt{tsc -b}
first). Tour ring misplaced after resize: it re-measures on resize and
route change automatically.

% =====================================================================
\chapter{Deployment}
% =====================================================================
The app is a static SPA bundle served as a Render \textbf{Web Service}:
\begin{table}[h]
\centering
\caption{Render service settings (repo \texttt{Devansh-567/baghewala-oil-field}).\label{tab:render}}
\begin{tabular}{lp{9.5cm}}
\toprule
\textbf{Field} & \textbf{Value} \\ \midrule
Language / Region & Node / Oregon (US West) \\
Branch & \texttt{main} \\
Build command & \texttt{npm install; npm run build} \\
Start command & \texttt{npm start} \;(runs \texttt{vite preview --host 0.0.0.0 --port \$\{PORT:-4173\}}) \\
Compute & Free tier (0.1 CPU / 512\,MB) suffices for demo traffic \\
\bottomrule
\end{tabular}
\end{table}
Render injects \texttt{\$PORT} (default 10000); binding \texttt{0.0.0.0} is
mandatory, and \texttt{vite preview} supplies SPA fallback so deep links such
as \texttt{/well/BGW-07/css} return \texttt{index.html} (verified: HTTP 200 on
\texttt{/} and \texttt{/field}). Local \texttt{npm start} on Windows/cmd needs
an explicit port since \texttt{\$\{PORT:-\dots\}} is sh syntax --- e.g.\
\texttt{npx vite preview --host 127.0.0.1 --port 4180}. Alternative: a Render
\textbf{Static Site} with publish directory \texttt{dist/} (no start command).

% =====================================================================
\chapter{Results and Validation}
% =====================================================================
\begin{table}[h]
\centering
\caption{Calibration checks (all satisfied by construction + smoke tests).\label{tab:valid}}
\begin{tabular}{lp{5cm}l}
\toprule
\textbf{Check} & \textbf{Expected} & \textbf{Result} \\ \midrule
$\mu$ anchor @50$^\circ$C & 8,000--15,000\,cP, API-ordered & By Eq.~\ref{eq:mu50} \\
$\mu(47^\circ$C) / $\mu(150^\circ$C) &
\textasciitilde13.5k\,cP / \textasciitilde12\,cP & Matches comment audit \\
Per-well rate & 15--45\,BOPD & Vogel + $F_p=0.42$ \\
Field scaling & $\approx$1,202\,BOPD @33 wells & Shown beside 12-pad total \\
Cycle SOR & 1--4, target $<3$ & 60-day flush window \\
Production build & \texttt{tsc -b} + Vite success & \textasciitilde900 modules \\
Serve smoke test & HTTP 200 on \texttt{/}, \texttt{/field} & Passed \\
\bottomrule
\end{tabular}
\end{table}
BGW-07 (showcase) solves to cycle 5, API 15.1, 680\,mD, pump 5.4\,SPM with
positive float margin and sub-0.55 Goodman at its recommended point; pushing
SPM to 6.6 inverts both --- the pump-off cliff evaluators can reproduce in
seconds.

% =====================================================================
\chapter{Limitations and Future Work}
% =====================================================================
\begin{enumerate}
  \item \textbf{No live SCADA yet:} providers are isolated (\texttt{src/data/*})
    behind stable shapes, so a real feed swaps in without touching pages; the
    drift-exclusion contract is already implemented.
  \item \textbf{Productivity factor $F_p=0.42$} is literature-typical, not
    history-matched to proprietary well files --- flagged openly in-app.
  \item \textbf{Lumped heat model:} Marx--Langenheim + soak efficiency, not a
    full thermal simulator (no CMG STARS coupling); multilayer and steam-override
    effects are future work.
  \item \textbf{Representative layout:} 12 pads compress the 52-well pattern;
    full-field rollout needs the operator's survey coordinates.
  \item \textbf{Roadmap:} history matching, per-well $F_p$ inversion, VFD
    closed-loop setpoint export, Hindi/English bilingual UI, PWA offline tile
    packs for field tablets.
\end{enumerate}

% =====================================================================
\chapter*{Conclusion}
\addcontentsline{toc}{chapter}{Conclusion}
The twin delivers what SIH26120 asked: CSS and SRP optimised \emph{together}
through one auditable chain, predictive heating/cooling/rate forecasts,
force-based rod-float detection, and steam/energy minimisation in rupees per
day --- presented on a real map, narrated by a guided tour, and traceable to
published data and textbook equations. It is a decision-support prototype with
an honest boundary: calibrated, not live --- and engineered so that going live
is a data-plumbing task, not a rewrite.

% =====================================================================
\begin{thebibliography}{99}
\addcontentsline{toc}{chapter}{References}
\bibitem{ML59} J.~W.~Marx and R.~H.~Langenheim, ``Reservoir heating by hot fluid
injection,'' \textit{Trans.\ AIME}, vol.~216, pp.~312--315, 1959.
\bibitem{BL66} T.~C.~Boberg and R.~B.~Lantz, ``Calculation of the production
rate of a thermally stimulated well,'' \textit{J.\ Pet.\ Technol.}, Oct.\ 1966
(SPE 1578-PA).
\bibitem{Vogel68} J.~V.~Vogel, ``Inflow performance relationships for
solution-gas drive wells,'' \textit{J.\ Pet.\ Technol.}, Jan.\ 1968.
\bibitem{API11L} American Petroleum Institute, \textit{RP 11L: Design
Calculations for Sucker-Rod Pumping Systems}, 5th ed.
\bibitem{API11E} American Petroleum Institute, \textit{Spec.\ 11E: Specification
for Pumping Units}.
\bibitem{API11BR} American Petroleum Institute, \textit{RP 11BR: Care and
Handling of Sucker Rods} (coupling SCF practice).
\bibitem{EJ92} T.~A.~Everitt and J.~W.~Jennings, ``An improved finite-difference
calculation of downhole dynamometer cards,'' SPE 18189, 1992.
\bibitem{Mehrotra} A.~K.~Mehrotra and W.~Y.~Svrcek, ``Viscosity of compressed
Athabasca bitumen,'' \textit{Can.\ J.\ Chem.\ Eng.}
\bibitem{Arps45} J.~J.~Arps, ``Analysis of decline curves,''
\textit{Trans.\ AIME}, vol.~160, 1945.
\bibitem{IAPWS} IAPWS, \textit{Revised Release on the IAPWS Formulation 1995
for the Thermodynamic Properties (IF97 steam tables)}.
\bibitem{OILRF} Oil India Ltd., ``Rajasthan Fields --- Baghewala,''
\url{https://www.oil-india.com/rajasthan-fields} (field data page).
\bibitem{SPE535203} Singh et al., ``Case study for enhancement of production of
heavy and highly viscous crude,'' SPE APOG 2023, Paper 535203.
\bibitem{DGH} Directorate General of Hydrocarbons, National Data Repository,
``Rajasthan Basin'' page.
\bibitem{SIH} Smart India Hackathon 2026, Problem Statement SIH26120
(Oil India Limited).
\bibitem{TOI26} Times of India / News18 / ET Energy, Apr.\ 2026, Oil India
record-production reports (1,202\,BOPD; CSS in 19 wells).
\bibitem{Bhat12} G.~M.~Bhat et al., ``Geology and hydrocarbon potential of
Neoproterozoic--Cambrian succession,'' \textit{Lyell Special Publications}
SP366 (Jodhpur Sandstone context).
\bibitem{React} Meta, React 18 documentation, \url{https://react.dev}.
\bibitem{Leaflet} Leaflet 1.9 + OpenStreetMap tile usage policy
(\url{https://operations.osmfoundation.org/policies/tiles/}).
\bibitem{Recharts} Recharts 2.x documentation, \url{https://recharts.org}.
\end{thebibliography}

% =====================================================================
\appendix
\chapter{Judge Q\&A (one-line answers)}
\begin{table}[h]
\centering
\begin{tabular}{p{5.5cm}p{9.5cm}}
\toprule
\textbf{Question} & \textbf{Answer} \\ \midrule
Real or synthetic data? & Calibrated model: published bands reproduced by
construction; layout representative; no live SCADA claimed. \\
Why 12 wells, not 52? & Demo compression; totals scale linearly to the record
(shown beside every total). \\
Novelty vs CSS software? & The coupling: one fillage state joins heating to
mechanics with INR-ranked optimisation. \\
How is rod float detected? & Downstroke force balance; margin in lbf
(900/250 thresholds). \\
What is SOR here? & CWE steam / oil in the 60-day flush window; target $<3$. \\
Can it take live SCADA? & Yes: isolated providers + drift-exclusion contract
already implemented. \\
Biggest limitation? & $F_p=0.42$ not field history-matched --- stated openly. \\
\bottomrule
\end{tabular}
\end{table}

\chapter{Glossary}
\begin{itemize}
  \item \textbf{CSS} --- Cyclic Steam Stimulation: inject $\rightarrow$ soak
    $\rightarrow$ produce in the same well.
  \item \textbf{SRP} --- Sucker-rod pump: beam unit + rod string + downhole
    plunger.
  \item \textbf{SOR} --- Steam-oil ratio; lower is better.
  \item \textbf{BOPD / BHT / VIT / PML} --- barrels of oil per day; bottomhole
    temperature; vacuum-insulated tubing; petroleum mining lease.
  \item \textbf{IPR / Goodman / SPM / CWE} --- inflow performance relationship;
    rod fatigue check; strokes per minute; cold-water equivalent.
  \item \textbf{Pareto knee} --- the operating point where extra steam cost
    starts outrunning production gain.
\end{itemize}

\chapter{Screenshot slots (Appendix C)}
Figures~\ref{fig:map}--\ref{fig:pareto} are placeholders. Save 1920$\times$1080
PNGs into \texttt{figures/} with the names below; recompile to embed them.
\begin{figure}[h]\centering
\fbox{\parbox{0.9\textwidth}{\centering\vspace{18pt}
\texttt{figures/field-overview.png}\\ \textit{Field Overview: OSM map + KPIs}\vspace{18pt}}}
\caption{Field overview (insert screenshot).}\label{fig:map}\end{figure}
\begin{figure}[h]\centering
\fbox{\parbox{0.9\textwidth}{\centering\vspace{18pt}
\texttt{figures/well-ipr.png}\\ \textit{Well twin: IPR + $\mu(T)$ charts}\vspace{18pt}}}
\caption{Well twin charts (insert screenshot).}\label{fig:ipr}\end{figure}
\begin{figure}[h]\centering
\fbox{\parbox{0.9\textwidth}{\centering\vspace{18pt}
\texttt{figures/css-sweeps.png}\\ \textit{CSS sweeps: steam vs soak economics}\vspace{18pt}}}
\caption{CSS optimisation (insert screenshot).}\label{fig:css}\end{figure}
\begin{figure}[h]\centering
\fbox{\parbox{0.9\textwidth}{\centering\vspace{18pt}
\texttt{figures/srp-cards.png}\\ \textit{Surface + pump dynacards at 6.6\,SPM}\vspace{18pt}}}
\caption{SRP cards (insert screenshot).}\label{fig:cards}\end{figure}
\begin{figure}[h]\centering
\fbox{\parbox{0.9\textwidth}{\centering\vspace{18pt}
\texttt{figures/optimizer-pareto.png}\\ \textit{Pareto frontier + top-8 table}\vspace{18pt}}}
\caption{Optimiser (insert screenshot).}\label{fig:pareto}\end{figure}

\end{document}
